\documentclass[10pt,a4paper]{amsart}
\usepackage{amsmath,amssymb,amsthm,graphicx}
% Minimal class-compatibility preamble for a reader excerpt of
%   Bhimarthi Ravinder, "Demazure Modules, Chari-Venkatesh Modules and Fusion Products",
%   SIGMA 10 (2014), 110, DOI 10.3842/SIGMA.2014.110.
%
% The author's pinned TeX file declares no theorem environments of its own: they are
% supplied by the SIGMA publisher class sigma.cls, which is not installed here.  The
% block below is that class's theorem-environment block (same environment names and
% styles), copied mechanically so that the retained author text compiles unchanged.
% No mathematical macro or statement is introduced by this file.
\newtheorem{theorem}{Theorem}
\newtheorem*{theorem*}{Theorem}
\newtheorem{lemma}{Lemma}
\newtheorem*{lemma*}{Lemma}
\newtheorem{corollary}{Corollary}
\newtheorem*{corollary*}{Corollary}
\newtheorem{proposition}{Proposition}
\newtheorem*{proposition*}{Proposition}
\newtheorem{conjecture}{Conjecture}
\newtheorem*{conjecture*}{Conjecture}
{\theoremstyle{definition} \newtheorem{definition}{Definition}
\newtheorem*{definition*}{Definition}
\newtheorem{example}{Example}
\newtheorem*{example*}{Example}
\newtheorem{remark}{Remark}
\newtheorem*{remark*}{Remark}
\newtheorem{note}{Note}
\newtheorem*{note*}{Note}
}
\usepackage{cleveref}
\crefname{originalreference}{original reference}{original references}
\Crefname{originalreference}{Original reference}{Original references}
\makeatletter
\newlabel{r1}{{r1}{source}{}{sourcepointer-1}{}}
\newlabel{r1@cref}{{[originalreference][0][]r1}{[1][0][]source}{}{}{}}
\makeatother
\crefname{originalreference}{original reference}{original references}
\Crefname{originalreference}{Original reference}{Original references}
\makeatletter
\newlabel{r2}{{r2}{source}{}{sourcepointer-2}{}}
\newlabel{r2@cref}{{[originalreference][0][]r2}{[1][0][]source}{}{}{}}
\makeatother
\crefname{originalreference}{original reference}{original references}
\Crefname{originalreference}{Original reference}{Original references}
\makeatletter
\newlabel{r3}{{r3}{source}{}{sourcepointer-3}{}}
\newlabel{r3@cref}{{[originalreference][0][]r3}{[1][0][]source}{}{}{}}
\makeatother
\crefname{originalreference}{original reference}{original references}
\Crefname{originalreference}{Original reference}{Original references}
\makeatletter
\newlabel{r4}{{r4}{source}{}{sourcepointer-4}{}}
\newlabel{r4@cref}{{[originalreference][0][]r4}{[1][0][]source}{}{}{}}
\makeatother
\crefname{originalreference}{original reference}{original references}
\Crefname{originalreference}{Original reference}{Original references}
\makeatletter
\newlabel{r5}{{r5}{source}{}{sourcepointer-5}{}}
\newlabel{r5@cref}{{[originalreference][0][]r5}{[1][0][]source}{}{}{}}
\makeatother
\crefname{originalreference}{original reference}{original references}
\Crefname{originalreference}{Original reference}{Original references}
\makeatletter
\newlabel{g}{{g}{source}{}{sourcepointer-6}{}}
\newlabel{g@cref}{{[originalreference][0][]g}{[1][0][]source}{}{}{}}
\makeatother
\crefname{originalreference}{original reference}{original references}
\Crefname{originalreference}{Original reference}{Original references}
\makeatletter
\newlabel{crucial}{{crucial}{source}{}{sourcepointer-7}{}}
\newlabel{crucial@cref}{{[originalreference][0][]crucial}{[1][0][]source}{}{}{}}
\makeatother
\title{Existence of the surjective map onto the kernel of the first comparison map between the modules attached to consecutive parameters}
\author{Bhimarthi Ravinder}
\date{doi:10.3842/SIGMA.2014.110; CC BY 4.0}
\begin{document}
\maketitle
\noindent\textit{Editorial note.} All delivered mathematical text is the author's own, copied from the pinned article source sources/sigma-2014-110-demazure-fusion/pinned-original-source.tex; no statement or proof step was rewritten, reordered or removed.\par\medskip
\noindent\textit{Editorial note.} Preamble provenance: the author's tex declares no theorem environments of its own (they come from the SIGMA publisher class); the staged sources/sigma-2014-110-demazure-fusion/retained-preamble.tex reproduces that class's theorem-environment block verbatim and is used as the preamble, with the wrapper class and base packages supplied by the assembler.\par\medskip
\noindent\textit{Editorial note.} Bibliography provenance: neither the counted statement nor its proof contains a citation (this lane rescans the whole span before writing the card), so the card carries no bibliography and no entry was reconstructed.\par\medskip
\noindent\textit{Editorial note.} Selection provenance: this card was opened under the file-based control contract (CONTROL.md), with the mechanical pre-card gate run on its statement and proof interval before the card was written.\par\medskip
\noindent\textit{Allowed standard prerequisites (editorial).}\begin{itemize}
\item The cyclic graded modules V\_\{i,k\} with their five families of defining relations, labelled r1 to r5 in the article, and the shift operation tau. The author's section 4.2, where the modules and their relations are set up before the counted statement; the five relations are referred to in the counted proof (source lines 641-648) as navigation-pointer stubs. Supplies the relations against which the map has to be checked, and the module on which it is defined by the value on one generator.
\item The two relations of the local Weyl module and of the Demazure module recorded in the lemmas labelled g and crucial. Both are invoked at the end of the counted proof (source line 649) as navigation-pointer stubs, each being counted separately in this lane. Identify the last two defining relations with statements already proved for the current algebra modules.
\end{itemize}\medskip
\ifdefined\hypertarget\hypertarget{sourcepointer-1}{}\fi
\noindent\textit{Source reference.} Original reference r1: sources/sigma-2014-110-demazure-fusion/pinned-original-source.tex, source lines 641--641. Navigation pointer only: the counted text refers to the item r1 of this same article, which is not part of this card. The pointer carries the item's own author label as its number; its content is not reproduced here.\par\medskip
\ifdefined\hypertarget\hypertarget{sourcepointer-2}{}\fi
\noindent\textit{Source reference.} Original reference r2: sources/sigma-2014-110-demazure-fusion/pinned-original-source.tex, source lines 647--647. Navigation pointer only: the counted text refers to the item r2 of this same article, which is not part of this card. The pointer carries the item's own author label as its number; its content is not reproduced here.\par\medskip
\ifdefined\hypertarget\hypertarget{sourcepointer-3}{}\fi
\noindent\textit{Source reference.} Original reference r3: sources/sigma-2014-110-demazure-fusion/pinned-original-source.tex, source lines 647--647. Navigation pointer only: the counted text refers to the item r3 of this same article, which is not part of this card. The pointer carries the item's own author label as its number; its content is not reproduced here.\par\medskip
\ifdefined\hypertarget\hypertarget{sourcepointer-4}{}\fi
\noindent\textit{Source reference.} Original reference r4: sources/sigma-2014-110-demazure-fusion/pinned-original-source.tex, source lines 648--648. Navigation pointer only: the counted text refers to the item r4 of this same article, which is not part of this card. The pointer carries the item's own author label as its number; its content is not reproduced here.\par\medskip
\ifdefined\hypertarget\hypertarget{sourcepointer-5}{}\fi
\noindent\textit{Source reference.} Original reference r5: sources/sigma-2014-110-demazure-fusion/pinned-original-source.tex, source lines 648--648. Navigation pointer only: the counted text refers to the item r5 of this same article, which is not part of this card. The pointer carries the item's own author label as its number; its content is not reproduced here.\par\medskip
\ifdefined\hypertarget\hypertarget{sourcepointer-6}{}\fi
\noindent\textit{Source reference.} Original reference g: sources/sigma-2014-110-demazure-fusion/pinned-original-source.tex, source lines 492--492. Navigation pointer only: the counted text refers to the item g of this same article, which is not part of this card. The pointer carries the item's own author label as its number; its content is not reproduced here.\par\medskip
\ifdefined\hypertarget\hypertarget{sourcepointer-7}{}\fi
\noindent\textit{Source reference.} Original reference crucial: sources/sigma-2014-110-demazure-fusion/pinned-original-source.tex, source lines 649--649. Navigation pointer only: the counted text refers to the item crucial of this same article, which is not part of this card. The pointer carries the item's own author label as its number; its content is not reproduced here.\par\medskip
\begin{proposition}
There exist a~surjective morphism of $\mathfrak{g}[t]$-modules $\phi^{-} : \tau_{2k+1-i} \mathbf{V}_{i,k}
\rightarrow \ker \phi^{+}$,
such that $\phi^{-}(\overline{v}_{i,k}) = \big(x^{-}_\theta \otimes t^{2k+1-i}\big) \overline{v}_{i,k+1}$.
\end{proposition}

\begin{proof}
We only need to show that $\phi^{-}(\overline{v}_{i,k})$ satisf\/ies the def\/ining relations of $\mathbf{V}_{i,k}$. We
start with the relation~\eqref{r1}.
First, for $\alpha=\theta$ it is clear.
Let $\alpha \in R^{+}\setminus\{\theta\}$; if $(\theta|\alpha)=0$ then also it is clear.
If $(\theta|\alpha)=1$ then $(\theta-\alpha)\in R^{+}\setminus\{\theta\}$ and $(\theta|\theta-\alpha)=1$, now it is
clear from the relations $(x^{-}_{\theta-\alpha}\otimes t^r)\overline{v}_{i,k+1}=0$ for all $r\geq (k+1)$ in
$\mathbf{V}_{i,k+1}$.
The relations~\eqref{r2},~\eqref{r3} are trivially satisf\/ied by $\phi^{-}(\overline{v}_{i,k})$.
Finally the last two relations~\eqref{r4},~\eqref{r5} are also satisf\/ied by $\phi^{-}(\overline{v}_{i,k})$; in fact
these are exactly the statements of Lemmas~\ref{g} and~\ref{crucial}.
\end{proof}

\end{document}
